\documentclass[11pt]{article}
\usepackage[a4paper,top=1.5cm,bottom=3.0cm,left=1.8cm,right=1.8cm,footskip=2.0cm]{geometry}
\usepackage{graphicx}
\usepackage[table]{xcolor}
\usepackage{ragged2e}
\usepackage{array}
\usepackage{fancyhdr}
\usepackage{fontspec}
\setmainfont{Noto Sans Gurmukhi}
\setsansfont{Noto Sans Gurmukhi}
\newfontfamily\latinfont{Carlito}
\renewcommand{\familydefault}{\sfdefault}
\definecolor{titleblue}{HTML}{0F2F4A}
\definecolor{textblack}{HTML}{1A1A1A}
\definecolor{rule}{HTML}{C9C4C0}
\definecolor{muted}{HTML}{6B635E}
\definecolor{rowalt}{HTML}{F3F1EF}
\setlength{\parindent}{0pt}\setlength{\parskip}{0.5em}

\setlength{\headheight}{0pt}\setlength{\headsep}{0pt}
\pagestyle{fancy}\fancyhf{}
\renewcommand{\headrulewidth}{0pt}\renewcommand{\footrulewidth}{0pt}
\fancyfoot[C]{%
  \begin{minipage}{\textwidth}
    {\color{rule}\hrule height 0.7pt}\vspace{4pt}
    \noindent
    \begin{minipage}[c]{0.66\textwidth}{\scriptsize\color{muted}{\latinfont Developed by the Water and Climate Lab, IIT Gandhinagar. Forecasts are driven by the IMD Extended Range Forecast System}.}\end{minipage}\hfill
    \begin{minipage}[c]{0.32\textwidth}\raggedleft
      
\includegraphics[height=0.62cm]{wcl.png}\hspace{6pt}%
      
\includegraphics[height=0.62cm]{iitgn.png}\hspace{6pt}%
      
\includegraphics[height=0.62cm]{imd.png}
    \end{minipage}\\[3pt]
    {\scriptsize\color{muted}© {\latinfont 2026 Water and Climate Lab} · {\latinfont Indian Institute of Technology Gandhinagar} · {\latinfont For research and demonstration purposes}.}
  \end{minipage}%
}

\begin{document}

\noindent
\begin{minipage}[c]{0.16\textwidth}
\includegraphics[width=\linewidth]{wcl.png}\end{minipage}\hfill
\begin{minipage}[c]{0.80\textwidth}\raggedleft
  {\fontsize{18}{21}\selectfont\bfseries\color{titleblue}{\latinfont Soil Moisture Forecast Outlook}}\\[2pt]
  {\normalsize\color{muted}{\latinfont Forecast initialized May 20, 2026}  ·  {\latinfont valid through Jun 19, 2026}}
\end{minipage}
\vspace{6pt}{\color{rule}\hrule height 1pt}\vspace{10pt}
{\large\bfseries\color{titleblue}{\latinfont Forecast maps (7 / 15 / 30 day)}}\\[4pt]
\noindent
\begin{minipage}[t]{0.323\textwidth}\centering{\small\bfseries\color{muted}{\latinfont 7-day}}\\[2pt]\includegraphics[width=\linewidth]{map7.png}\end{minipage}\hfill
\begin{minipage}[t]{0.323\textwidth}\centering{\small\bfseries\color{muted}{\latinfont 15-day}}\\[2pt]\includegraphics[width=\linewidth]{map15.png}\end{minipage}\hfill
\begin{minipage}[t]{0.323\textwidth}\centering{\small\bfseries\color{muted}{\latinfont 30-day}}\\[2pt]\includegraphics[width=\linewidth]{map30.png}\end{minipage}
\\[3pt]{\small\color{muted}{\latinfont Forecast soil moisture for the 7-, 15- and 30-day horizons, initialized May 20, 2026}.}\par
\vspace{10pt}
{\large\bfseries\color{titleblue}{\latinfont Outlook}}\\[2pt]
{\justifying\normalsize\color{textblack}ਰਾਸ਼ਟਰੀ ਪੱਧਰ 'ਤੇ ਉਮੀਦ ਕੀਤਾ ਜਾਂਦਾ ਹੈ ਕਿ ਮਿੱਟੀ ਦੀ ਨਮੀ ਦਾ ਔਸਤ ਲਗਭਗ {\latinfont 0.20 v/v (7}-ਦਿਨ), {\latinfont 0.20 v/v (15}-ਦਿਨ) ਅਤੇ {\latinfont 0.21 v/v (30}-ਦਿਨ) ਹੈ।\par}
\vspace{8pt}
{\large\bfseries\color{titleblue}{\latinfont Regional Outlook}}\\[3pt]
{\normalsize\color{muted}{\latinfont Expected soil moisture (v/v)}}\\[4pt]
\renewcommand{\arraystretch}{1.3}
\begin{tabular}{@{}p{0.30\textwidth} >{\raggedleft\arraybackslash}p{0.16\textwidth} >{\raggedleft\arraybackslash}p{0.16\textwidth} >{\raggedleft\arraybackslash}p{0.16\textwidth}@{}}
\rowcolor{titleblue}
\textcolor{white}{\textbf{ {\latinfont Region} }} & \textcolor{white}{\textbf{ {\latinfont 7-day} }} & \textcolor{white}{\textbf{ {\latinfont 15-day} }} & \textcolor{white}{\textbf{ {\latinfont 30-day} }}\\
\rowcolor{rowalt} {\latinfont North} & {\latinfont 0.20} & {\latinfont 0.20} & {\latinfont 0.20}\\
 {\latinfont Northwest} & {\latinfont 0.15} & {\latinfont 0.15} & {\latinfont 0.15}\\
\rowcolor{rowalt} {\latinfont Northeast} & {\latinfont 0.33} & {\latinfont 0.34} & {\latinfont 0.35}\\
 {\latinfont East} & {\latinfont 0.17} & {\latinfont 0.17} & {\latinfont 0.18}\\
\rowcolor{rowalt} {\latinfont Central} & {\latinfont 0.17} & {\latinfont 0.17} & {\latinfont 0.20}\\
 {\latinfont West} & {\latinfont 0.21} & {\latinfont 0.22} & {\latinfont 0.25}\\
\rowcolor{rowalt} {\latinfont South} & {\latinfont 0.21} & {\latinfont 0.22} & {\latinfont 0.23}\\
\end{tabular}
\end{document}
